% GENERATED FILE. Edit canonical lessons, then run npm run generate:books.
% canonical-source-hash: fnv1a64:d3729a228d1f53a9
% canonical-lessons: IT-C266-dovrei, IT-C266-spiegare, IT-C266-ripasso-spiegazione

\chapter{Dovrei, Spiegare --- Advice, and the Shape of an Explanation}
\label{ch:it-advice-and-explanation}
\hlchaptermodality{266}

\begin{chapteropening}
\textbf{By the end of this chapter:} \emph{I can advise with dovrei and dovresti, and explain myself in four moves: an opinion, a cause, a consequence and advice.}

Complete IT-C266-ripasso-spiegazione: answer a why-question with all four moves of an explanation.
\end{chapteropening}

\section[dovrei]{dovrei — what you ought to do about it}

\label{lesson:IT-C266-dovrei}



\begin{quote}
Ask for a coffee politely. (\emph{Vorrei un caffè.}) Then run a reason forward: \textquotedblleft{}I'm tired, so I'm staying at home.\textquotedblright{} (\emph{Sono stanco, quindi resto a casa.})

An explanation usually ends on a third move: what somebody ought to do. That move has its own form.
\end{quote}

\subsection*{You'll want to know: dovrei}
You have \textbf{dovere}, \textquotedblleft{}to have to, must\textquotedblright{}. Soften it and it becomes advice:

\begin{quote}
\textbf{Dovrei} studiare. — I should study. \textbf{Dovresti} dormire. — You should sleep. \emph{(to a friend)} \textbf{Dovrebbe} riposare. — He or she should rest; and, to someone you call \emph{Lei}, \textquotedblleft{}you should rest\textquotedblright{}.
\end{quote}

A bare infinitive follows, with no \emph{di} and no \emph{che}, exactly as after \emph{vorrei}. And listen to the ending: \textbf{dovrei}, \textbf{vorrei}. It is the same \emph{-ei}. Both are the \textbf{condizionale}, the \textquotedblleft{}would\textquotedblright{} form of Italian, which you have been using since chapter eighty-six without its name. Chapter two hundred and seventy-two builds it properly. For now, two forms of it are yours.

\emph{Devo studiare} is \textquotedblleft{}I must study\textquotedblright{}: no way out. \emph{Dovrei studiare} is \textquotedblleft{}I should\textquotedblright{}: I know I ought to, and I have not promised. That softness is the register an explanation usually wants.

\begin{cousinweb}
\textbf{dovere} is inherited from Latin \emph{dēbēre}, \textquotedblleft{}to owe\textquotedblright{} (EtymDB 2.1), and \emph{dēbēre} is a compound: \emph{dē}, \textquotedblleft{}from\textquotedblright{}, plus \emph{habēre}, \textquotedblleft{}to have\textquotedblright{} (Lewis \& Short). To owe is to \textbf{have something from someone}: you are holding what is theirs. You met the noun in chapter two hundred and fifty-three: \textbf{il debito}, \textquotedblleft{}a debt\textquotedblright{}. English took the same Latin word as \textbf{debt}, \textbf{due} and \textbf{duty}.
\end{cousinweb}

\subsection*{Your turn}
Say these aloud:

\begin{itemize}
\raggedright
  \item \emph{dovrei, dovresti, dovrebbe}
  \item the hard one, then the soft one — \emph{Devo studiare.} then \emph{Dovrei studiare.}
  \item advice to a tired friend — \emph{Sei stanco, quindi dovresti dormire.}
  \item \emph{dovrei} and \emph{vorrei} together — the same \emph{-ei}
\end{itemize}

\begin{tcolorbox}[breakable,colback=teal!4,colframe=teal!35!black,title={Before you move on}]
What is the difference between \emph{devo} and \emph{dovrei}? (\emph{\textbf{Devo}} is \textquotedblleft{}I must\textquotedblright{}; \emph{\textbf{dovrei}} is \textquotedblleft{}I should\textquotedblright{} — softer.) Which form you already knew has the same ending? (\emph{\textbf{Vorrei.}}) What are the two Latin pieces of \emph{dēbēre}? (\emph{\textbf{Dē}} + \emph{\textbf{habēre}} — to have from someone.)
\end{tcolorbox}
\section[spiegare]{spiegare — to unfold}

\label{lesson:IT-C266-spiegare}



\begin{quote}
Three pieces are in your hands: \emph{secondo me} for an opinion, \emph{perché} and \emph{quindi} for the two arrows, \emph{dovrei} for advice. This lesson puts them in a row, under the verb that names the act.
\end{quote}

\begin{grammarlens}[title={Grammar Lens: the shape of an explanation}]
You met \textbf{spiegare}, \textquotedblleft{}to explain\textquotedblright{}, in chapter one hundred and forty-one. It is a regular \emph{-are} verb with one spelling catch: \textbf{spiego}, \textbf{spieghi}, \textbf{spiega}. The \emph{h} in \emph{spieghi} keeps the \emph{g} hard before the \emph{i}, the same trick the \emph{h} plays in \emph{che}.

Here is a whole explanation, every word of it taught:

\begin{quote}
— Perché resti a casa? — \textbf{Secondo me} piove. E ho il raffreddore, \textbf{perché} ieri faceva freddo. \textbf{Quindi} resto a casa. \textbf{Dovrei} dormire.
\end{quote}

\begin{quote}
\emph{— Why are you staying at home?} \emph{— If you ask me, it's going to rain. And I've got a cold, because it was cold yesterday. So I'm staying at home. I ought to sleep.}
\end{quote}

Four moves, and a word for each: an \textbf{opinion} (\emph{secondo me}), a \textbf{cause} pointing back (\emph{perché}), a \textbf{consequence} pointing forward (\emph{quindi}), and \textbf{advice} (\emph{dovrei}). With both arrows, you can build it cause first or claim first.
\end{grammarlens}

\begin{cousinweb}
\textbf{piegare}, \textquotedblleft{}to fold\textquotedblright{}, which you met in chapter one hundred and eighty, is inherited from Latin \emph{plicāre}, \textquotedblleft{}to fold\textquotedblright{} (EtymDB 2.1). \textbf{spiegare} is the same verb with Latin \emph{ex-}, \textquotedblleft{}out\textquotedblright{}, worn down to an \emph{s-}: Latin \emph{explicāre}, \textquotedblleft{}to unfold, to spread out\textquotedblright{} (Lewis \& Short, under \emph{explico}). EtymDB 2.1 has no row for \emph{spiegare}, so that pairing is the usual reading rather than a cited one.

\textbf{To explain is to unfold}: something was folded shut, and you open it out. English kept the Latin in \textbf{explicit}, and \textbf{complicated} is the opposite, folded together. \textbf{explain} is another verb, Latin \emph{explānāre}, \textquotedblleft{}to make plain\textquotedblright{} (EtymDB 2.1).
\end{cousinweb}

\subsection*{Your turn}
Say these aloud:

\begin{itemize}
\raggedright
  \item \emph{spiego, spieghi, spiega} — keep the \emph{g} hard
  \item the four moves — \emph{Secondo me piove, perché non c'è il sole, quindi resto a casa. Dovrei leggere.}
  \item the same, consequence first — \emph{Resto a casa perché piove.}
  \item \emph{piegare}, \emph{spiegare}, then English \emph{explicit}
\end{itemize}

\begin{tcolorbox}[breakable,colback=teal!4,colframe=teal!35!black,title={Before you move on}]
Name the four moves of an explanation and a word for each. (Opinion \emph{\textbf{secondo me}}, cause \emph{\textbf{perché}}, consequence \emph{\textbf{quindi}}, advice \emph{\textbf{dovrei}}.) Why is there an \emph{h} in \emph{spieghi}? (To keep the \emph{\textbf{g}} \textbf{hard}.) What is \emph{spiegare}, taken apart? (\emph{\textbf{Piegare}} with \emph{ex-}: \textbf{to unfold}.)
\end{tcolorbox}
\section[(review)]{Review — explain yourself}

\label{lesson:IT-C266-ripasso-spiegazione}



\begin{quote}
Say \textquotedblleft{}you should sleep\textquotedblright{} to a friend. (\emph{Dovresti dormire.}) Which English word for owing shares its Latin? (\textbf{Debt}, or \textbf{due}, or \textbf{duty}.)
\end{quote}

\begin{grammarlens}[title={What you've built}]
Two chapters, four moves. You can say what you think, why, what follows from it, and what ought to be done. And because chapter two hundred and sixty-three gave you the scene, your reasons can reach into the past:

\begin{quote}
— Perché non sei andato al mercato? — Faceva freddo ed ero stanco, quindi sono restato a casa. Secondo me dovrei uscire di più.
\end{quote}

The reason is scene, \emph{faceva} and \emph{ero}; the consequence is an event, \emph{sono restato}; the advice is about now. Three times in one answer.
\end{grammarlens}

\subsection*{Your turn}
Say these aloud:

\begin{itemize}
\raggedright
  \item the question — \emph{Perché non sei andato al mercato?}
  \item the answer, all four moves — \emph{Faceva freddo ed ero stanco, quindi sono restato a casa. Secondo me dovrei uscire di più.}
  \item a different reason — \emph{Pioveva, così ho studiato.}
  \item answer \emph{Dovresti studiare?} with \emph{Credo di sì.}
  \item an explanation of your own: an opinion, a \emph{perché}, a \emph{quindi}, a \emph{dovrei}
\end{itemize}

\begin{tcolorbox}[breakable,colback=teal!4,colframe=teal!35!black,title={Before you move on}]
In \emph{faceva freddo, quindi sono restato a casa}, which half is the scene and which the event? (\emph{\textbf{Faceva freddo}} is the scene; \emph{\textbf{sono restato}} is the event.) What does \emph{spiegare} do, literally? (It \textbf{unfolds}.) And what is softer than \emph{devo}? (\emph{\textbf{Dovrei.}})
\end{tcolorbox}
